% from ega2-1-fr.tex
\oldpage[II]{101}

Il suffit de démontrer (i), (ii) et (iii) (\hyperref[I.3.5.1-fr]{I, 3.5.1}).
Dans chacun de ces trois cas, la vérification de la condition a) de
(\hyperref[II.5.4.1-fr]{5.4.1}) découle de résultats antérieurs
(\hyperref[I.5.5.1-fr]{I, 5.5.1} et \hyperref[I.6.3.4-fr]{6.3.4})~; reste à
vérifier la condition b). C'est immédiat dans le cas (i), car si $X\to Y$
est une immersion fermée, il en est de même de
$X\times_YY'\to Y\times_YY'=Y'$ (\hyperref[I.4.3.2-fr]{I, 4.3.2} et
\hyperref[II.3.3.3-fr]{3.3.3}). Pour démontrer (ii), considérons deux
morphismes propres $X\to Y$, $Y\to Z$, et un morphisme $Z'\to Z$. On peut
écrire $X\times_ZZ'=X\times_Y(Y\times_ZZ')$
(\hyperref[I.3.3.9.1-fr]{I, 3.3.9.1}), et par suite la projection
$X\times_ZZ'\to Z'$ se factorise en
$X\times_Y(Y\times_ZZ')\to Y\times_ZZ'\to Z'$. Compte tenu de la remarque
du début, (ii) résulte donc de ce que le composé de deux morphismes fermés
est fermé. Enfin, pour tout morphisme $S'\to S$, $X_{(S')}$ s'identifie à
$X\times_YY_{(S')}$ (\hyperref[I.3.3.11-fr]{I, 3.3.11})~; pour tout
morphisme $Z\to Y_{(S')}$, on peut écrire
\[
  X_{(S')}\times_{Y_{(S')}}Z
  =(X\times_YY_{(S')})\times_{Y_{(S')}}Z=X\times_YZ~;
\]
$X\times_YZ\to Z$ est fermé par hypothèse, donc cela prouve (iii).

\begin{corollary}[5.4.3]
\label{II.5.4.3-fr}
Soient $f:X\to Y$, $g:Y\to Z$ deux morphismes tels que $g\circ f$ soit
propre.
\begin{enumerate}
  \item[\rm{(i)}] Si $g$ est séparé, $f$ est propre.
  \item[\rm{(ii)}] Si $g$ est séparé et de type fini, et si $f$ est
  surjectif, $g$ est propre.
\end{enumerate}
\end{corollary}

(i) découle de (\hyperref[II.5.4.2-fr]{5.4.2}) par le procédé général
(\hyperref[I.5.5.12-fr]{I, 5.5.12}). Pour démontrer (ii), il n'y a à
vérifier que la condition b) de la déf. (\hyperref[II.5.4.1-fr]{5.4.1}).
Pour tout morphisme $Z'\to Z$, le diagramme
\[
  \xymatrix{
    X\times_ZZ'\ar[r]^{f\times1_{Z'}}\ar[dr]_p &
    Y\times_ZZ'\ar[d]^{p'}\\
    & Z'
  }
\]
(où $p$ et $p'$ sont les projections) est commutatif
(\hyperref[I.3.2.1-fr]{I, 3.2.1})~; en outre, $f\times1_{Z'}$ est surjectif
puisqu'il en est ainsi de $f$ (\hyperref[I.3.5.2-fr]{I, 3.5.2}), et $p$ est
un morphisme fermé par hypothèse. Toute partie fermée $F$ de
$Y\times_ZZ'$ est alors l'image par $f\times1_{Z'}$ d'une partie fermée $E$
de $X\times_ZZ'$, donc $p'(F)=p(E)$ est fermé dans $Z'$ par hypothèse, d'où
le corollaire.

\begin{corollary}[5.4.4]
\label{II.5.4.4-fr}
Si $X$ est un préschéma propre au-dessus de $Y$, et $\mathcal{S}$ une
$\mathcal{O}_Y$-Algèbre quasi-cohérente, tout $Y$-morphisme
$f:X\to\operatorname{Proj}(\mathcal{S})$ est propre (et a fortiori fermé).
\end{corollary}

En effet, le morphisme structural
$p:\operatorname{Proj}(\mathcal{S})\to Y$ est séparé, et $p\circ f$ est
propre par hypothèse.

\begin{corollary}[5.4.5]
\label{II.5.4.5-fr}
Soit $f:X\to Y$ un morphisme séparé de type fini. Soient
$(X_i)_{1\leq i\leq n}$ (resp. $(Y_i)_{1\leq i\leq n}$) une famille finie
de sous-préschémas fermés de $X$ (resp. $Y$), $j_i$ (resp. $h_i$)
l'injection canonique $X_i\to X$ (resp. $Y_i\to Y$). Supposons que l'espace
sous-jacent $X$ soit réunion des $X_i$, et que pour tout $i$, il y ait un
morphisme $f_i:X_i\to Y_i$ tel que le diagramme
\[
  \xymatrix{
    X_i\ar[r]^{f_i}\ar[d]_{j_i} & Y_i\ar[d]^{h_i}\\
    X\ar[r]_f & Y
  }
\]
soit commutatif. Alors, pour que $f$ soit propre, il faut et il suffit que
chacun des $f_i$ le soit.
\end{corollary}

